\documentclass[10pt,letterpaper]{article}
\usepackage{technical-design-spec}
\usepackage{graphicx}
\usepackage{longtable,booktabs,array}
\usepackage{listings}
\setDocTitle{AI Agent Harness Course Architecture Description}
\setDocSubtitle{C4 Models, UML Diagrams and ISO/IEC/IEEE 42010:2022 Assessment}
\setDocVersion{0.1}
\setDocStatus{Draft -- partial conformance}
\setDocOwner{AI Agent Engineering}
\setDocAuthor{Technical Architecture}
\setDocDate{24 November 2026}
\setDocSummary{Evidence-based reverse engineering of the supplied agent harness course snapshot. The architecture description is a static self-assessment, not evidence of production deployment or certification.}
\begin{document}
\makedoctitle
\tableofcontents
\listoffigures
\clearpage
\section{Executive summary and identification}
AD-AHC-001, revision 0.1, 24 November 2026. Entity of interest: the collection of agent-harness teaching implementations in richmondalake/agent\_harness\_course. The supplied snapshot label is 8a5edab282632443; full commit identity, branch and checkout freshness are not verified. The text dump contains 585 file sections (216 in Part 1, 207 in Part 2, 160 in Part 3, two at root). This description concerns architecture, not functional certification. Its purpose is to communicate boundaries, responsibilities, behavior, data and risk to learners, developers and reviewers. Environment: browser-facing Appbook examples, Python runtimes, configurable local and live services, optional databases, MCP workspace tools and remote model/sandbox providers. Scope excludes production claims, live provider calls, private data, credential exposure and application changes.

Part 1 demonstrates retail planning, packaged versus custom harnesses, and advanced supplier-workflow, ontology, research and evaluation examples. Part 2 moves to personal productivity with workspace integration and durable research. Part 3 extends the productivity baseline with cost accounting, trajectories, explicit learning gates and more elaborate research. These are distinct educational implementations, not three tiers of a single deployed system. Prior six-document study is a companion, not an independently verified review.

\section{Method and evidence policy}
Repository sections were inventoried; architecture-significant code symbols and configuration locations were selected for static review. Evidence identifiers refer to file-relative line ranges where available, not the dump's line numbers. Direct code observations are distinguished from inferred grouping, proposed controls and unknown runtime behavior. A module name establishes existence, not correctness or isolation. See the evidence CSV and the coverage register. No application scripts, benchmarks, network calls or live tests were executed. No licensed reference or source dump is redistributed. Reviewers should recheck the actual checkout and all file-relative line anchors before accepting this draft.
\section{Stakeholders, perspectives, concerns and aspects}
Roles are representative, not verified identities or approvals. The JSON model records many-to-many associations. Concern questions define the intended coverage; missing implementation evidence remains a known issue rather than an invented feature.
\begin{longtable}{@{}>{\raggedright\arraybackslash}p{0.15\linewidth}>{\raggedright\arraybackslash}p{0.75\linewidth}@{}}
\toprule
ID & Stakeholder\\\midrule\endhead
ST01 & Assistant user\\\addlinespace
ST02 & Course instructor/learner\\\addlinespace
ST03 & Developer/architect\\\addlinespace
ST04 & Operator/data steward\\\addlinespace
ST05 & Security/QA reviewer\\\addlinespace
ST06 & Course owner/maintainer\\\addlinespace
\bottomrule\end{longtable}
\begin{longtable}{@{}>{\raggedright\arraybackslash}p{0.15\linewidth}>{\raggedright\arraybackslash}p{0.75\linewidth}@{}}
\toprule
ID & Perspective\\\midrule\endhead
P01 & System/educational\\\addlinespace
P02 & Engineering\\\addlinespace
P03 & Information\\\addlinespace
P04 & Operational\\\addlinespace
P05 & Security\\\addlinespace
P06 & Evaluation\\\addlinespace
\bottomrule\end{longtable}
\begin{longtable}{@{}>{\raggedright\arraybackslash}p{0.15\linewidth}>{\raggedright\arraybackslash}p{0.75\linewidth}@{}}
\toprule
ID & Concern\\\midrule\endhead
C01 & What is in scope, and how do course variants relate?\\\addlinespace
C02 & How are application responsibilities and boundaries allocated?\\\addlinespace
C03 & Which real code constructs implement the design?\\\addlinespace
C04 & How does inference become a bounded execution loop?\\\addlinespace
C05 & How are workflows paused, resumed and recovered?\\\addlinespace
C06 & How is memory retrieved, persisted and governed?\\\addlinespace
C07 & How are external tools and model providers integrated?\\\addlinespace
C08 & What execution and storage environments are configured?\\\addlinespace
C09 & How are costs, failures and traces observed?\\\addlinespace
C10 & Who authorizes consequential operations?\\\addlinespace
C11 & How are untrusted input and sandbox exposure managed?\\\addlinespace
C12 & How are learning, evaluation and architecture evolution governed?\\\addlinespace
\bottomrule\end{longtable}
\begin{longtable}{@{}>{\raggedright\arraybackslash}p{0.15\linewidth}>{\raggedright\arraybackslash}p{0.75\linewidth}@{}}
\toprule
ID & Aspect\\\midrule\endhead
A01 & Structural\\\addlinespace
A02 & Behavioral\\\addlinespace
A03 & Information\\\addlinespace
A04 & Physical/operational\\\addlinespace
A05 & Security\\\addlinespace
A06 & Evolutionary\\\addlinespace
\bottomrule\end{longtable}
\section{Viewpoint and model-kind specifications}
Each viewpoint is version 0.1 and governs the associated V-VPxx view. Each view includes narrative plus the registered diagrams assigned to it, where available. Model kinds and view methods are defined in the JSON model. C4 diagrams use the project's pinned C4-PlantUML style; UML diagrams use corresponding native syntax. Interpret boxes by type, labels by direction, and guards as conditional transitions. A source identifier establishes traceability, not completeness. View construction extracts source entities and call sites; interpretation compares scope and relationships with the source, and validation checks references, syntax, rendering and visual clarity.
\begin{longtable}{@{}>{\raggedright\arraybackslash}p{0.22\linewidth}>{\raggedright\arraybackslash}p{0.22\linewidth}>{\raggedright\arraybackslash}p{0.47\linewidth}@{}}
\toprule
Viewpoint & Concerns / aspects & Method\\\midrule\endhead
VP01 Context and landscape & C01 / A01 & Identify boundaries and actors; compare separate baselines, not a single deployment.\\\addlinespace
VP02 Functional and logical & C02 / A01 A02 & Inspect entry points and responsibility allocation; distinguish logical module from process.\\\addlinespace
VP03 Implementation & C03 / A01 & Extract declared classes and dependencies; avoid invented interfaces.\\\addlinespace
VP04 Runtime and harness & C04 C05 / A02 & Trace graph edges and call sites; compare proposal, execution and persistence.\\\addlinespace
VP05 Information and memory & C06 / A03 & Trace memory sources, stores and promotion; identify unverified retention.\\\addlinespace
VP06 Integration & C07 / A01 A02 & Trace provider, MCP and tool contracts and failure translation.\\\addlinespace
VP07 Deployment and operations & C08 C09 / A04 & Map logical to configured execution; do not infer observed deployment.\\\addlinespace
VP08 Security & C10 C11 / A05 A02 & Trace authority, trust and approval; separate risks from exploit findings.\\\addlinespace
VP09 Evaluation and evolution & C12 / A02 A06 & Trace evaluation and skill versions; compare architecture alternatives and limitations.\\\addlinespace
\bottomrule\end{longtable}
\begin{longtable}{@{}>{\raggedright\arraybackslash}p{0.26\linewidth}>{\raggedright\arraybackslash}p{0.65\linewidth}@{}}
\toprule
Model kind & Notation and interpretation\\\midrule\endhead
C4 context & C4 persons/systems and directional relationship labels; landscape uses multiple systems\\\addlinespace
C4 container & C4 system boundary and applications/data stores\\\addlinespace
C4 component & C4 container boundary and logical components\\\addlinespace
C4 deployment & C4 nodes and container instances; configured environment only\\\addlinespace
UML sequence & UML participants and ordered messages; alt denotes branches\\\addlinespace
UML state & UML states, transitions and guards; conceptual labels noted\\\addlinespace
UML package & UML packages and dependency arrows\\\addlinespace
UML class & UML classes and declared operations\\\addlinespace
UML activity & UML actions and activity flow\\\addlinespace
Narrative & Non-model view component: headings, evidence IDs, uncertainty labels\\\addlinespace
\bottomrule\end{longtable}
\section{Structural and functional analysis}
The packaged MemoRizz assistant and custom ERPA graph are alternative teaching examples (E01, E04). ERPA's offline graph is explicit context assembly, model decision, tool dispatch and persistence. Its live graph and Oracle factories are separate (E02--E03); parity cannot be assumed. The cache sits at the boundary before context assembly (E01, E29), so a cache hit returns without recall or new graph persistence. The trade-off is reduced model work at the expense of stale policy or memory risk; inspect keys and invalidation before reuse across users. The custom ERPA graph groups tools by intent rather than exposing an unrestricted arbitrary executor. Advanced workflow, Total Recall, research and MetaHarness have distinct entry points (E06--E08, E28, E33--E34). The landscape intentionally leaves their unverified internal dependencies absent.

The PPA Part 2 graph, approval module and scheduler allocate orchestration, effect governance and scheduled work (E10, E13--E14). Workspace MCP exposes mail, calendar and notes; its provider adapter wraps failures as a structured result (E11). That is an integration boundary, not blanket permission. The Part 3 graph reuses the productivity problem and adds recorded context size, per-call ledger, trajectory capture and evaluation machinery (E16--E20). The Appbook container grouping in C4 L2 is an inferred logical teaching runtime; browser and server can be decomposed where independent execution is established. Data storage is logical, not a claim of a single shared database.

\section{Runtime architecture and authority}
Part 3 build\_graph compiles named nodes with a supplied checkpointer (E16). Context assembly reads standing memory and up to six recalled records of selected types, records context layers and prepares the request. The model adapter produces a response and tool proposals. The graph records usage before dispatch. Dispatch validates arguments, rejects unavailable model-facing tools, evaluates tier and refusal, and prevents tool execution after the configured iteration budget; replay forces effectful calls into approval. This is a bounded proposal/execution design, not proof that every other call path enforces the same policy.

For approval-tier calls the harness assesses risk, drafts an action and marks the run awaiting approval. The graph interrupts; resumption reads the durable action record. Only APPROVED is executed, with arguments validated again; a still-drafted action is rejected. The model is not the authority and the resume payload is not the gate (E16). The execution path should be tested for binding to the approved arguments, principal, expiry, atomic state and replay behavior. The graph returns to inference after effects; persist writes user/assistant messages and marks the run completed. Checkpoint durability is distinct from exactly-once external side effects.

The Part 2 research graph has outline review, gather/read, organize, fan-out writing, review, assembly, publication review and close (E15). Part 3 includes tally, reflect and budget helpers (E21, E35) plus evidence deduplication and excerpt selection (E22). The researcher subgraph (E27) is separate from the outer paper workflow. These are workflow and evaluator-like patterns; no comparative accuracy or throughput was measured. Part 1 workflow API exposes start/resume/decision/inspection (E06); the actual restart trace is untested.

\section{Information and memory}
ERPA context contains memory, skills, tool descriptions and semantic information; explicit user preferences and working-turn information are persisted (E01). Local and Oracle providers differ (E03, E34). PPA memory has a provider contract and local implementation (E24); Part 3 uses standing and typed recalled records in the prompt (E16). The logical information classes are requests, replies, tool arguments/results, memory records, checkpoint state, action audit rows, skills/versions, usage events, trajectory events and evidence sources. Model-generated or externally retrieved content is untrusted; the audit row is an authority for a specific approved effect, not authority for unrelated operations. Physical tables and retention are not exhaustively mapped. Proposed production controls: principal-bound namespaces, explicit provenance, retention/deletion enforcement, protected promotion and isolation tests.

Learning captures repeated trajectories and evaluates stored skill versions through real-tool, regression, held-out, approval-contract and replay checks. A temporary live pointer is restored after the trial; the gate records PASSED or FAILED. Promotion requires PASSED, while a separate activation operation can move the pointer for rollback (E17--E18). Cost is reported beside gate results, not a pass criterion. A production design should enforce activation permissions and prevent unapproved candidate versions from being activated via an alternate route.
\section{Integration, deployment and operations}
The external model provider supplies inference via an adapter; the application owns validation and effects. FastMCP provides stateless HTTP workspace tools in Part 2; the analogous Part 3 edge is inferred, and protocol/authentication must be tested (E11, E16, E30). Oracle factories and local/Oracle storage selection exist (E03, E23), but configuration names do not establish deployed resources. The Part 1 Dockerfile and Compose configuration provide packaging evidence (E25--E26); the deployment view is a conceptual mapping, not a verified live topology. Inspect the actual manifests and resolve runtime boundaries for each profile before operational approval.

The ledger records per-call cost and usage (E19), and trajectories record planning, close and failure events (E20). Do not equate estimates or scripted usage with billing. Observability should correlate run, thread, model invocation, tool proposal, decision, effect and checkpoint without exporting sensitive payloads. Availability, SLO, recovery time, retry behavior and capacity remain unmeasured.
\section{Security architecture assessment}
Trust boundaries include browser-to-application, application-to-model, model output-to-dispatch, dispatch-to-MCP/provider, application-to-persistence, human decision-to-audit, and host-to-sandbox. Protect credentials, workspace records, business data, durable state, authorization records and skill versions. STRIDE framing: spoofing of user/decision, tampering with approved arguments or memory, repudiation of effects, disclosure through logs and provider messages, denial through unbounded work, escalation through tools or skill activation. The following are risk hypotheses based on code, not confirmed vulnerabilities or exploit results.
\begin{longtable}{@{}>{\raggedright\arraybackslash}p{0.18\linewidth}>{\raggedright\arraybackslash}p{0.39\linewidth}>{\raggedright\arraybackslash}p{0.35\linewidth}@{}}
\toprule
Risk / priority & Scenario and existing control & Proposed verification\\\midrule\endhead
R01 High (E16 E32) & Approval replay or argument mutation between drafting and execution. Audit state checked, arguments revalidated; atomic binding and idempotency unverified. & Bind canonical arguments, principal, action digest and expiry; test denial/replay.\\\addlinespace
R02 High (E11 E12) & Indirect prompt injection via workspace documents or model outputs. Heuristic wrapping and recipient-risk helpers; no proof against adversarial content. & Treat retrieved text as data; capability policy independent of model; adversarial fixtures.\\\addlinespace
R03 High (E05) & Local subprocess treated as security sandbox; host-side callable executes before isolation. Remote provider requests no egress and closes session; test profile explicitly unproven. & Use least-privilege remote isolation; restrict host callables; test filesystem, egress and timeout.\\\addlinespace
R04 High (E17 E18) & Unsafe skill activation or poisoned learned procedure. Promotion checks PASSED; separate activate path can move pointer to another version. & Authorize activation separately, enforce eligible rollback versions; provenance and held-out tests.\\\addlinespace
R05 High (E01 E29) & Cache hit bypasses current memory and policy context. Cache placed ahead of graph; exact scoping and invalidation not verified. & Include tenant, policy, model and context versions in key; cross-tenant and stale-policy tests.\\\addlinespace
R06 High (E24 E16) & Memory poisoning, cross-user retrieval or stale preferences. Typed recall and bounded result count; isolation and source trust unverified. & Namespace records, preserve provenance, enforce promotion policy, test poisoning and deletion.\\\addlinespace
R07 Medium (E16 E19 E20) & Telemetry captures sensitive requests and tool data. Usage ledger and trajectories exist; redaction and retention unverified. & Structured redaction, restricted access, retention and deletion tests.\\\addlinespace
R08 High (E15 E21 E16) & Duplicate writes after restart or interruption. Checkpoint graphs and approval rows; no executed restart proof in this review. & Idempotency keys, transactional outbox, crash-injection tests around side effects.\\\addlinespace
R09 Medium (E11) & MCP bearer token and error translation can expose information or insufficiently scope permissions. Token configured; full auth and provider error handling not evaluated. & Per-tool scopes, fail closed if missing token, redact errors and test direct server access.\\\addlinespace
R10 Medium (E21 E22 E35) & Research evidence poisoning, runaway cost or publication without provenance. Budget/reflection and evidence functions exist; performance/effectiveness unmeasured. & Citation verification, budget fail-closed tests, source isolation, publication authorization.\\\addlinespace
\bottomrule\end{longtable}
\paragraph{Sandbox qualification.} The Part 1 advanced executor resolves and validates registered arguments, runs a trusted host callable to derive an expected result, then selects remote E2B or local subprocess, checks remote no-egress configuration, forwards an empty environment, verifies output equality and closes the provider in a finally block (E05). Host recomputation is not isolation; local fallback explicitly does not prove egress blocking. Do not execute arbitrary model-supplied code with host authority.
\section{Quality attributes and evaluation}
The following acceptance scenarios are proposed, not measured. For each, stimulus is supplied by a reviewer in a test environment, the target is the named component, and the response measure is an observable test oracle.
\begin{longtable}{@{}>{\raggedright\arraybackslash}p{0.18\linewidth}>{\raggedright\arraybackslash}p{0.38\linewidth}>{\raggedright\arraybackslash}p{0.36\linewidth}@{}}
\toprule
Attribute & Stimulus and target & Proposed response measure\\\midrule\endhead
Security & Unapproved effect or replay to dispatcher & Zero external writes; decision logged\\\addlinespace
Reliability & Interrupt after effect before checkpoint & No duplicate effect on restart; correlated audit record\\\addlinespace
Privacy & Cross-principal recall and cache request & Zero foreign records or answers returned\\\addlinespace
Performance & Repeated workload against cache/graph & Record p50/p95 and model calls; define SLA before acceptance\\\addlinespace
Maintainability & Replace model or persistence provider & Contract tests pass without domain changes\\\addlinespace
Cost & Budget-exhausting research task & Stops or requests explicit grant; ledger reconciles\\\addlinespace
Evaluation & Poisoned skill or unknown tool & Gate fails and live pointer unchanged\\\addlinespace
Observability & Injected provider failure & Run/action/trace linkage with redacted error\\\addlinespace
\bottomrule\end{longtable}
\section{Architecture evolution and decisions}
Part 1 demonstrates a sequential custom loop plus packaged alternative; Part 2 introduces workspace effect governance and a durable paper workflow; Part 3 strengthens observability and learning evaluation and extends research with budgeting/reflection. Dependency separation, provider abstraction, graph routing, human-in-the-loop and fan-out are visible patterns. Hexagonal ports/adapters as a formal architecture are proposed rather than asserted for the entire repository. Parallel execution, supervisor delegation and systematic production load testing are not established across all variants.
\begin{longtable}{@{}>{\raggedright\arraybackslash}p{0.23\linewidth}>{\raggedright\arraybackslash}p{0.24\linewidth}>{\raggedright\arraybackslash}p{0.45\linewidth}@{}}
\toprule
Decision & Observed status and evidence & Inferred rationale / trade-off\\\midrule\endhead
ADR01 Separate teaching baselines & Observed choice; motivation inferred; E01 E04 E10 E16 & Isolate variants rather than deploy them as one system. Trade-off: duplicated modules versus easier comparison.\\\addlinespace
ADR02 Cache before graph & Observed choice; motivation inferred; E01 E29 & Low-latency cache hit avoids work; trades freshness and policy context for speed. Alternative: cache only after authority context.\\\addlinespace
ADR03 Approval persisted separately from model & Observed choice; motivation inferred; E16 E32 & Human-controlled audit row decouples resume from permission; requires atomicity and idempotency.\\\addlinespace
ADR04 Checkpoint graph & Observed choice; motivation inferred; E16 E15 E21 & Durability improves restartability but introduces versioning, replay and state migration concerns.\\\addlinespace
ADR05 Skill gate and version pointer & Observed choice; motivation inferred; E17 E18 & Evaluation before promote and rollback support; separate activation path needs governance.\\\addlinespace
ADR06 Pluggable storage and model adapters & Observed choice; motivation inferred; E03 E23 E24 E30 & Local teaching profile and live dependencies; interface drift and semantic parity must be tested.\\\addlinespace
ADR07 Host-recomputed sandbox result & Observed choice; motivation inferred; E05 & Output integrity check is valuable but host execution must be trusted; alternative: pure verification only.\\\addlinespace
ADR08 Cost reporting separate from learning gate & Observed choice; motivation inferred; E17 E19 & Cost visibility without confusing affordability with correctness; alternative: separate explicit budget policy.\\\addlinespace
\bottomrule\end{longtable}
\subsection{Proposed future-state reference}
Keep the domain/application core provider-neutral, with typed model, memory, tools, authorization, checkpoint and telemetry ports. A policy decision point receives principal, action, resource, argument digest and context; a policy enforcement point rechecks immediately before effect execution. Use a durable idempotency key and outbox, quarantine untrusted retrieval, validate skill versions in an isolated trial, and expose explicit operator controls. Independent adapters provide MCP, model SDK, database and remote sandbox implementations. This proposal is not the observed course implementation.
\section{Architecture views and diagram register}
All diagrams have editable PlantUML sources beside this document. Where the repository's canonical renderer has generated PNGs, they are embedded; otherwise a source locator is printed rather than substituting an unverified figure. A diagram's governing view, model kind, evidence and interpretation follow. The diagrams are draft; successful rendering and semantic review remain gates.
\newcommand{\archfigure}[3]{\IfFileExists{diagrams/png/#1.png}{\begin{figure}[htbp]\centering\includegraphics[width=\linewidth,height=.66\textheight,keepaspectratio]{diagrams/png/#1.png}\caption{#2}\label{fig:#1}\end{figure}}{\par\noindent\textit{Render pending. Source: }\path{diagrams/#1.puml}\par}\noindent #3\par}
\subsection{C4-L1-01 -- Part 1 system context}
\textbf{View:} V-VP01; \textbf{model kind:} C4 context; \textbf{evidence:} E01, E02.\par
\archfigure{c4-l1-01}{Part 1 system context}{Retail planning; offline mode is scripted and live model use is configuration dependent.}
\clearpage

\subsection{C4-L1-02 -- Part 2 system context}
\textbf{View:} V-VP01; \textbf{model kind:} C4 context; \textbf{evidence:} E10, E11.\par
\archfigure{c4-l1-02}{Part 2 system context}{Personal productivity and workspace integration; deployment is not observed.}
\clearpage

\subsection{C4-L1-03 -- Part 3 system context}
\textbf{View:} V-VP01; \textbf{model kind:} C4 context; \textbf{evidence:} E16, E30.\par
\archfigure{c4-l1-03}{Part 3 system context}{Learning-enabled PPA; learning is inside the system, not an autonomous external service.}
\clearpage

\subsection{C4-LS-001 -- Course implementation landscape}
\textbf{View:} V-VP01; \textbf{model kind:} C4 context; \textbf{evidence:} E01, E04, E06, E07, E08, E15, E16, E21, E28.\par
\archfigure{c4-ls-001}{Course implementation landscape}{Conceptual landscape of independent examples; arrows denote study, not runtime calls.}
\clearpage

\subsection{C4-L2-001 -- Part 1 logical containers}
\textbf{View:} V-VP02; \textbf{model kind:} C4 container; \textbf{evidence:} E01, E02, E03.\par
\archfigure{c4-l2-001}{Part 1 logical containers}{Logical Appbook grouping and persistence are inferred; split browser client if independently runnable and confirm actual deployment. MCP counterpart for Part 3 inferred.}
\clearpage

\subsection{C4-L2-002 -- Part 2 logical containers}
\textbf{View:} V-VP02; \textbf{model kind:} C4 container; \textbf{evidence:} E10, E11.\par
\archfigure{c4-l2-002}{Part 2 logical containers}{Logical Appbook grouping and persistence are inferred; split browser client if independently runnable and confirm actual deployment. MCP counterpart for Part 3 inferred.}
\clearpage

\subsection{C4-L2-003 -- Part 3 logical containers}
\textbf{View:} V-VP02; \textbf{model kind:} C4 container; \textbf{evidence:} E16, E23, E30.\par
\archfigure{c4-l2-003}{Part 3 logical containers}{Logical Appbook grouping and persistence are inferred; split browser client if independently runnable and confirm actual deployment. MCP counterpart for Part 3 inferred.}
\clearpage

\subsection{C4-L3-001 -- Part 1 Appbook components}
\textbf{View:} V-VP02; \textbf{model kind:} C4 component; \textbf{evidence:} E01, E03, E09, E29.\par
\archfigure{c4-l3-001}{Part 1 Appbook components}{Approval service is shown as a separate capability; no universal invocation edge is claimed.}
\clearpage

\subsection{C4-L3-002 -- Part 2 Appbook components}
\textbf{View:} V-VP02; \textbf{model kind:} C4 component; \textbf{evidence:} E10, E12, E13, E14.\par
\archfigure{c4-l3-002}{Part 2 Appbook components}{Part 2 dependency directions are a working hypothesis requiring call-site verification.}
\clearpage

\subsection{C4-L3-003 -- Part 3 Appbook components}
\textbf{View:} V-VP02; \textbf{model kind:} C4 component; \textbf{evidence:} E16, E17, E18, E19, E20, E24, E30, E31, E32.\par
\archfigure{c4-l3-003}{Part 3 Appbook components}{Learning and runtime components share the teaching app but learning is not shown as automatically triggered by each turn.}
\clearpage

\subsection{UML-SQ-001 -- ERPA cache hit and miss}
\textbf{View:} V-VP04; \textbf{model kind:} UML sequence; \textbf{evidence:} E01, E29.\par
\archfigure{uml-sq-001}{ERPA cache hit and miss}{Offline run method; cache hit bypasses context and persistence in the graph. Scope/invalidation require tests.}
\clearpage

\subsection{UML-SQ-002 -- Part 3 bounded agent loop}
\textbf{View:} V-VP04; \textbf{model kind:} UML sequence; \textbf{evidence:} E16, E30.\par
\archfigure{uml-sq-002}{Part 3 bounded agent loop}{Simplified iteration; gated branch expanded in UML-SQ-003.}
\clearpage

\subsection{UML-SQ-003 -- Approval interrupt and audit authority}
\textbf{View:} V-VP08; \textbf{model kind:} UML sequence; \textbf{evidence:} E16, E32.\par
\archfigure{uml-sq-003}{Approval interrupt and audit authority}{The audit row, not the resume payload, is the observed execution gate. Authentication, binding and exactly-once behavior remain unverified.}
\clearpage

\subsection{UML-SQ-004 -- Memory context assembly}
\textbf{View:} V-VP05; \textbf{model kind:} UML sequence; \textbf{evidence:} E16, E24.\par
\archfigure{uml-sq-004}{Memory context assembly}{Shows observed read path only; promotion, retention and isolation require separate evidence.}
\clearpage

\subsection{UML-SQ-005 -- Workspace MCP integration}
\textbf{View:} V-VP06; \textbf{model kind:} UML sequence; \textbf{evidence:} E11, E16, E31.\par
\archfigure{uml-sq-005}{Workspace MCP integration}{Part 2 server failure wrapper verified; Part 3 integration edge inferred, not an authenticated trace.}
\clearpage

\subsection{UML-SQ-006 -- Skill trial and promotion}
\textbf{View:} V-VP09; \textbf{model kind:} UML sequence; \textbf{evidence:} E17, E18.\par
\archfigure{uml-sq-006}{Skill trial and promotion}{The distinct activate/rollback path is not equivalent to promote; validate authorization and permitted states.}
\clearpage

\subsection{UML-SQ-007 -- Run observability and cost}
\textbf{View:} V-VP07; \textbf{model kind:} UML sequence; \textbf{evidence:} E16, E19, E20, E30.\par
\archfigure{uml-sq-007}{Run observability and cost}{Ledger reports estimated/billed fields; no independent cost or latency measurement performed.}
\clearpage

\subsection{UML-SM-001 -- Agent run lifecycle}
\textbf{View:} V-VP04; \textbf{model kind:} UML state; \textbf{evidence:} E16.\par
\archfigure{uml-sm-001}{Agent run lifecycle}{Graph nodes and statuses observed. Crash/error state not added without checking all callers.}
\clearpage

\subsection{UML-SM-002 -- Approval lifecycle}
\textbf{View:} V-VP08; \textbf{model kind:} UML state; \textbf{evidence:} E16, E32.\par
\archfigure{uml-sm-002}{Approval lifecycle}{Executed is a conceptual terminal stage, not necessarily a literal database state.}
\clearpage

\subsection{UML-SM-003 -- Skill version lifecycle}
\textbf{View:} V-VP09; \textbf{model kind:} UML state; \textbf{evidence:} E17, E18.\par
\archfigure{uml-sm-003}{Skill version lifecycle}{Conceptual lifecycle; activation/rollback is a separate operation not exhaustively modeled.}
\clearpage

\subsection{UML-PK-001 -- Part 3 harness source packages}
\textbf{View:} V-VP03; \textbf{model kind:} UML package; \textbf{evidence:} E16, E17, E18, E19, E20, E24, E30, E31, E32.\par
\archfigure{uml-pk-001}{Part 3 harness source packages}{Actual module names; diagram edges reflect observed calls.}
\clearpage

\subsection{C4-L4-001 -- ERPA graph code-level class}
\textbf{View:} V-VP03; \textbf{model kind:} UML class; \textbf{evidence:} E01.\par
\archfigure{c4-l4-001}{ERPA graph code-level class}{Actual class and operations; no fabricated inheritance.}
\clearpage

\subsection{C4-L4-002 -- Memory provider code-level view}
\textbf{View:} V-VP03; \textbf{model kind:} UML class; \textbf{evidence:} E24.\par
\archfigure{c4-l4-002}{Memory provider code-level view}{Declared class names; inheritance and signature omitted until exact class declaration is checked.}
\clearpage

\subsection{UML-AC-001 -- Part 2 paper workflow}
\textbf{View:} V-VP04; \textbf{model kind:} UML activity; \textbf{evidence:} E15.\par
\archfigure{uml-ac-001}{Part 2 paper workflow}{Simplified successful path through named graph nodes; review rejection and restart omitted from this view.}
\clearpage

\subsection{UML-AC-002 -- Part 3 research workflow}
\textbf{View:} V-VP04; \textbf{model kind:} UML activity; \textbf{evidence:} E21, E22, E27, E35.\par
\archfigure{uml-ac-002}{Part 3 research workflow}{Budget and reflection added to paper workflow; no measured quality or cost claim.}
\clearpage

\subsection{SEC-AU-001 -- Part 1 sandbox execution boundary}
\textbf{View:} V-VP08; \textbf{model kind:} UML sequence; \textbf{evidence:} E05.\par
\archfigure{sec-au-001}{Part 1 sandbox execution boundary}{The registered callable executes on the host before sandbox execution. Only registered deterministic functions should enter this path; verify provider isolation independently.}
\clearpage

\subsection{C4-DP-001 -- Part 1 configured deployment boundary}
\textbf{View:} V-VP07; \textbf{model kind:} C4 deployment; \textbf{evidence:} E25, E26, E03.\par
\archfigure{c4-dp-001}{Part 1 configured deployment boundary}{Deployment mapping is inferred from Docker/Compose presence and Oracle factory; no claim about exact topology or a running instance.}
\clearpage
\section{Correspondences and consistency}
The architecture JSON records the participating AD, identifiers, and governing method for each relationship. COR01 checks stakeholder association, COR02 coverage of concerns, perspectives and aspects, COR03 view/component membership, and COR04 evidence and containment. The validator checks structural integrity and source presence, not semantic truth. Known issues: Part 3 MCP interaction edge inferred; Part 2 component dependencies inferred; Appbook grouping not a fully validated process decomposition; physical data and deployment mapping incomplete; other runtime scenarios not traced. Diagrams and source catalog require render and manual review. No consistency or conformance claim follows from a successful syntax check.
\section{ISO/IEC/IEEE 42010:2022 conformance}
Clause 4 distinguishes AD, ADF, ADL, viewpoint and model-kind claims. This draft assesses the AD case under Clause 6, incorporating referenced Clause 8 provisions. It makes no separate framework or language claim. No normative requirement is waived to obtain conformance. Status values are provisional document-inspection findings: satisfied, partially satisfied, not assessed and not applicable (separate claim only). The matrix records requirements separately, with verification evidence or gaps. Independent certification is not claimed.
\begin{longtable}{@{}>{\raggedright\arraybackslash}p{0.13\linewidth}>{\raggedright\arraybackslash}p{0.34\linewidth}>{\raggedright\arraybackslash}p{0.45\linewidth}@{}}
\toprule
Clause & Assessment item & Result and evidence/gap\\\midrule\endhead
6.1 & Entity and environment & Satisfied: Scope section\\\addlinespace
6.1 & Purpose & Satisfied: Scope section\\\addlinespace
6.1 & Identification metadata & Satisfied: Document header and baseline\\\addlinespace
6.2 & Relevant stakeholders & Partially satisfied: Representative roles; actual ownership unconfirmed\\\addlinespace
6.2 & Earlier evaluation recommendations considered & Partially satisfied: Companion study referenced; complete previous reviews unknown\\\addlinespace
6.2 & Resource constraints documented & Satisfied: Limitations\\\addlinespace
6.2 & Explain nonconformances & Satisfied: Assessment and gaps\\\addlinespace
6.3 & Identify perspectives & Satisfied: Model catalog\\\addlinespace
6.3 & Link stakeholders to perspectives & Satisfied: Model catalog\\\addlinespace
6.3 & Enumerate perspective concerns & Satisfied: Model catalog\\\addlinespace
6.4 & Identify concerns & Satisfied: Concern register\\\addlinespace
6.4 & Associate stakeholder and concern & Satisfied: Model and COR01\\\addlinespace
6.5 & Identify aspects & Satisfied: Model catalog\\\addlinespace
6.5 & Associate concerns and aspects & Satisfied: Model and COR02\\\addlinespace
6.6 & Include viewpoints & Satisfied: Viewpoint catalog\\\addlinespace
6.6 & Version viewpoints & Satisfied: Model version\\\addlinespace
6.6 & Viewpoint specifications meet clause 8 & Partially satisfied: Render and semantic review pending\\\addlinespace
6.6 & Frame every concern & Satisfied: COR02 structure\\\addlinespace
6.6 & Link perspectives and viewpoints & Satisfied: Model catalog\\\addlinespace
6.7 & View per viewpoint & Satisfied: View catalog\\\addlinespace
6.7 & Version every view & Satisfied: Model catalog\\\addlinespace
6.7 & Perspectives addressed & Partially satisfied: Actual adequacy not independently reviewed\\\addlinespace
6.7 & Concerns addressed & Partially satisfied: Scenario coverage incomplete\\\addlinespace
6.7 & Aspects addressed & Partially satisfied: Deployment verification outstanding\\\addlinespace
6.7 & View follows conventions & Not assessed: PlantUML rendering and manual review pending\\\addlinespace
6.7(a) & View identifying information & Satisfied: Model and diagram register\\\addlinespace
6.7(b) & Governing viewpoint & Satisfied: Model and diagram register\\\addlinespace
6.7(c) & View components address framed concerns & Partially satisfied: Source-based views; missing scenario coverage\\\addlinespace
6.7(d) & Known view issues & Satisfied: Catalog and limitations\\\addlinespace
6.8 & Compose views from view components & Satisfied: Model\\\addlinespace
6.8 & Version components & Satisfied: Model\\\addlinespace
6.8 & Identify model kind and conventions & Partially satisfied: Semantics and rendering unverified\\\addlinespace
6.8 & Legend for non-model content & Satisfied: Narrative legend\\\addlinespace
6.9.1 & Known inconsistencies recorded & Satisfied: Limitations and open issues\\\addlinespace
6.9.2 & Correspondence list & Satisfied: Model catalog\\\addlinespace
6.9.2 & Participating AD elements identified & Satisfied: Model links\\\addlinespace
6.9.2 & Governing correspondence methods identified & Satisfied: Model links\\\addlinespace
6.9.2 & Participating ADs identified & Satisfied: AD-AHC-001 on links\\\addlinespace
6.9.3 & List applicable correspondence methods & Satisfied: Model catalog\\\addlinespace
6.9.3 & Record method satisfaction or violations & Partially satisfied: Structural validator supplied; execution pending\\\addlinespace
6.9.3 & Include methods applying to AD & Satisfied: Model catalog\\\addlinespace
6.10.1 & Essential decisions recorded & Partially satisfied: Selected decisions; exhaustive history unavailable\\\addlinespace
6.10.2 & Rationale for essential decisions & Partially satisfied: Observed choices, rationale inferred\\\addlinespace
8.1(a) & Viewpoint perspectives & Satisfied: Model\\\addlinespace
8.1(b) & Viewpoint concerns & Satisfied: Model\\\addlinespace
8.1(c) & Viewpoint aspects & Satisfied: Model\\\addlinespace
8.1(d) & Stakeholders holding concerns & Satisfied: Model\\\addlinespace
8.1(e) & Model kinds and legends & Partially satisfied: Catalog uses notation; rendering pending\\\addlinespace
8.1(f) & Correspondence methods & Satisfied: Model\\\addlinespace
8.1(g) & Viewpoint sources & Satisfied: Model\\\addlinespace
8.1 & Legend interpretation guidance & Partially satisfied: Visual inspection pending\\\addlinespace
8.2(a) & Model notation conventions & Satisfied: Model kinds\\\addlinespace
8.2(b) & View and correspondence methods & Satisfied: Model kinds\\\addlinespace
8.2(c) & Version identification & Satisfied: Model kinds\\\addlinespace
8.2(d) & Model kind sources & Satisfied: Model kinds\\\addlinespace
8.3 & Necessary view methods defined & Partially satisfied: Construction method supplied; full analysis methods pending\\\addlinespace
7.1 & Separate ADF conformance case & Not applicable: No ADF conformance claim\\\addlinespace
7.2 & Separate ADL conformance case & Not applicable: No ADL conformance claim\\\addlinespace
\bottomrule\end{longtable}
\textbf{Conformance statement: not established.} Applicable requirements are partially assessed; unresolved verification and architecture coverage prevent a complete claim.
\appendix
\section{Coverage and unresolved work}
\begin{longtable}{@{}>{\raggedright\arraybackslash}p{0.34\linewidth}>{\raggedright\arraybackslash}p{0.58\linewidth}@{}}
\toprule
Scenario & Status\\\midrule\endhead
Assistant response & UML-SQ-001/002\\\addlinespace
Tool proposal/validation & UML-SQ-002/003\\\addlinespace
MCP & UML-SQ-005 (inferred path)\\\addlinespace
Memory retrieval & UML-SQ-004\\\addlinespace
Semantic search & Partial: E03; dedicated sequence absent\\\addlinespace
Cache & UML-SQ-001\\\addlinespace
Approval & UML-SQ-003 / UML-SM-002\\\addlinespace
Sandbox & SEC-AU-001\\\addlinespace
Workflow checkpoints & Partial: E15 E16 E21; restart sequence absent\\\addlinespace
Workflow recovery & Not assessed; no runtime restart test\\\addlinespace
Scheduled tasks & Partial: E14; dedicated lifecycle absent\\\addlinespace
Deep research & UML-AC-001/002\\\addlinespace
Evaluation/feedback & Partial: E08 E17\\\addlinespace
Memory promotion & Not fully traced\\\addlinespace
Skill creation/reuse & UML-SQ-006 / UML-SM-003 (draft/reuse not fully traced)\\\addlinespace
Deployment & C4-DP-001 inferred; other environments not modeled\\\addlinespace
Total Recall & C4-LS-001; internals not fully traced\\\addlinespace
MetaHarness & C4-LS-001; internals not fully traced\\\addlinespace
Packaged harness & C4-LS-001; internals not fully traced\\\addlinespace
\bottomrule\end{longtable}
\section{Evidence register}
File-relative paths, symbols and line anchors are preserved in the accompanying CSV. This summary gives the interpretation of each source; the snapshot is the baseline, not a deployed system.
\begin{longtable}{@{}>{\raggedright\arraybackslash}p{0.1\linewidth}>{\raggedright\arraybackslash}p{0.32\linewidth}>{\raggedright\arraybackslash}p{0.5\linewidth}@{}}
\toprule
ID & Symbol & Observation\\\midrule\endhead
E01 & ERPAStateGraph.run, assemble\_context, persist (64-183) & Offline cache short-circuits context; misses recall, decide, dispatch and persist; explicit preference writes and offline checkpoint.\\\addlinespace
E02 & compile\_live\_graph, LiveOracleGraph, invoke\_with\_cache (73-336) & Separate live LangGraph implementation; Oracle adapters and cache boundary.\\\addlinespace
E03 & build\_oamp, build\_langchain\_components, build\_checkpointer (125-270) & Oracle-backed live memory, retrieval and checkpoint factory.\\\addlinespace
E04 & MemAgent (51-100) & Packaged-harness baseline; distinct assistant entry point.\\\addlinespace
E05 & SandboxedToolExecutor.execute (311-421) & Registry schema validation, remote egress guard, empty forwarded env, host recomputation and finally-close; local subprocess is test profile.\\\addlinespace
E06 & start\_run, resume\_run, decide\_run, inspect\_run (145-175) & Workflow API exposes start, resume, decision and inspection.\\\addlinespace
E07 & TotalRecallState, TotalRecallHarness (26-641) & Total Recall implementation has its own state and harness; ontology-oriented workflow.\\\addlinespace
E08 & MetaHarnessCourse (71 onward) & MetaHarness evaluation/demo entry point; not evidence of production performance.\\\addlinespace
E09 & approve\_action, execute\_approved\_action (39-100) & Action approval state transition and separate execution.\\\addlinespace
E10 & build\_graph, Agent (185-374) & PPA graph and run facade; Part 2 baseline distinct from Part 3.\\\addlinespace
E11 & guarded, FastMCP, authorised (35-199) & Workspace MCP mail/calendar/notes capabilities and token-based gateway; provider error translated into result.\\\addlinespace
E12 & detect\_injection, wrap\_untrusted, recipient\_risk (51-99) & Heuristic content handling and recipient risk functions, not proof of injection resistance.\\\addlinespace
E13 & assess, draft, decide, executed (109-238) & Risk assessment, draft, decision and execution audit lifecycle.\\\addlinespace
E14 & Scheduler (74 onward) & Scheduled work component.\\\addlinespace
E15 & build\_graph, durable\_graph, start\_paper, resume\_paper (212-283) & Durable paper workflow, outline and publication review, fan-out writing.\\\addlinespace
E16 & build\_graph (212-408) & Recall/context, model, ledger, bounded dispatch, typed validation, replay gating, interrupt, audit-row authorization and persistence.\\\addlinespace
E17 & gate\_skill, promote, activate, retire (469-535) & Five checks and temporary trial; promotion checks PASSED; separate activation path merits authorization review.\\\addlinespace
E18 & activate, retire, trial, disclose (182-318) & Versioned live pointer, trial restoration and progressive disclosure.\\\addlinespace
E19 & record, calls, summary (83-117) & Per-call usage/cost ledger; scripted and live runs must not be conflated.\\\addlinespace
E20 & open\_turn, plan, close, fail (32-82) & Trajectory lifecycle captures turn evidence.\\\addlinespace
E21 & tally, reflect, review, build\_graph, durable\_graph (137-340) & Budget/reflect/review nodes and durable graph.\\\addlinespace
E22 & new\_results, keep\_source, choose\_excerpts (60-156) & Evidence deduplication and excerpt selection.\\\addlinespace
E23 & db\_path, checkpoint\_path, initialize (180-334) & Local/Oracle storage selection and initialization.\\\addlinespace
E24 & MemoryProvider, LocalMemoryProvider, build\_provider (34-171) & Memory provider contract, local implementation and expiry helper.\\\addlinespace
E25 & file (entire file) & Container packaging exists for Part 1; actual deployed topology unverified.\\\addlinespace
E26 & file (entire file) & Compose configuration exists; runtime deployment not observed.\\\addlinespace
E27 & researcher\_graph, research\_section (172-188) & Research agent subgraph separate from outer paper workflow.\\\addlinespace
E28 & start\_session, approve\_session, run\_session (187-235) & Research session creation, approval and execution routes.\\\addlinespace
E29 & BoundarySemanticCache (18 onward) & Cache implementation; scope and invalidation require targeted testing.\\\addlinespace
E30 & respond, structured, usage\_of (60-142) & Model response/structured-output and usage extraction adapter.\\\addlinespace
E31 & ToolContext, Args and tool registry (26 onward) & Typed tool argument classes; dispatch validation at E16.\\\addlinespace
E32 & assess, draft, decide, executed (109-238) & Approval API and audit transitions used by E16.\\\addlinespace
E33 & OracleHarnessRunStore, OracleApprovalStore (79-334) & Dedicated persistence classes for advanced runs and approvals.\\\addlinespace
E34 & InMemoryAgentMemoryManager, OracleAgentMemoryManager (42-424) & Alternative memory implementations.\\\addlinespace
E35 & usage, typical, estimate, grant (15-36) & Research budget calculations and grant helper.\\\addlinespace
\bottomrule\end{longtable}
\section{Diagram review and build}
Inspect every rendered diagram for title, scope, audience, element name/type/responsibility, readable relationship direction/label, explained notation and legend, consistent aliases and evidence. Confirm all depicted relationships against source or mark inferred; record deviations. Run the repository's PlantUML lint and renderer before building the PDF, then inspect figure cropping and reference resolution. Run the adjacent structural validator from this directory. A missing PNG is deliberately shown as render pending, not a successful figure. No generated binaries are checked in. Existing CI policy is unchanged.
\section{References}
\begin{itemize}
\item ISO/IEC/IEEE 42010:2022, \emph{Software, systems and enterprise -- Architecture description}, Clauses 4, 6, 8 and Annex B. Licensed supplied reference; not redistributed.
\item Simon Brown, \emph{The C4 Model: Visualizing Software Architecture}, O'Reilly, 2026, Chapters 2--10 and diagram review appendix.
\item C4 Model, \url{https://c4model.com/}; OMG UML, \url{https://www.omg.org/spec/UML/}; PlantUML, \url{https://plantuml.com/}; C4-PlantUML, \url{https://github.com/plantuml-stdlib/C4-PlantUML}.
\item Supplied course snapshot, label 8a5edab282632443; repository-specific evidence register accompanies this description.
\end{itemize}
\end{document}
